(1) 无限长均匀细弦质量线密度 \(\mu>0\)、弹性张力参数 \(T>0\)，具有Kelvin–Voigt应变率阻力系数 \(\eta>0\)（与 \(T\) 同类乘以时间量纲）。小横位移满足横向单位长度内力 \(T u_x+\eta u_{xt}\) 的空间梯度。推导波动方程及复色散关系，采用 \(e^{i(kx-\omega t)}\) 的频率符号约定。
(2) 求两支复频率并分类波数 \(0<|k|<k_c\)、\(|k|=k_c\) 与 \(|k|>k_c\) 的欠阻尼、临界、过阻尼情况。求长波相速度与幅度衰减率的首项，说明为何不能把实群速度推广为高 \(k\) 的传播速度。
(3) 对周期边界或空间无穷远衰减的任意实位移场，从波动方程严格推出机械能
\(E=\frac12\int[\mu u_t^2+Tu_x^2]\,dx\) 的耗散恒等式，辨别耗散机制是 \(u_{xt}\) 而非 \(u_t\)。
(4) 对固定 \(\eta,T,\mu\) 的短波极限推导两支纯衰减时间尺度的标度，并核对 \(\eta\to0^+\) 极限、特征波数的量纲以及在微观长度尺度为何应放弃局域连续介质近似。
